\chapter{Hansen's Disease}

\section*{Definition and Overview}
Hansen's disease, also called leprosy, is a chronic bacterial infection caused by \textit{Mycobacterium leprae}, which mainly affects the skin, peripheral nerves, and the lining of the nose. It is a curable infection, and modern multidrug therapy is highly effective. Hansen's disease is not highly contagious: it spreads through prolonged, close contact with untreated infected people, and around 95% of people have natural immunity. Untreated, it causes damage to the nerves, leading to loss of sensation, weakness, and disability. Early diagnosis and treatment prevent disability and cure the infection. Hansen's disease is rare in Australia, with a small number of new cases each year, mostly in migrants and in Aboriginal communities in remote northern Australia.

\section*{Epidemiology}
Hansen's disease is rare in Australia: fewer than 10 new cases are diagnosed per year, mostly among people who lived in or visited countries where the disease is more common, and among Aboriginal and Torres Strait Islander people in northern regions. Globally, it remains a significant public health issue, with around 200,000 new cases reported each year, mostly in India, Brazil, and parts of Africa and Southeast Asia. The disease affects people of all ages, and the risk of disability increases with delay in diagnosis. With treatment, the disease is curable and transmission stops.

\section*{Aetiology and Pathophysiology}
Hansen's disease is caused by \textit{Mycobacterium leprae}, a slow-growing bacterium.

\begin{itemize}
    \item \textbf{The organism:} \textit{Mycobacterium leprae} grows very slowly, with an incubation period of two to five years or longer
    \item \textbf{Transmission:} through prolonged, close contact with untreated infected people, likely via respiratory droplets and contact with nasal secretions; it is not spread by casual contact
    \item \textbf{Immune response:} the type of disease depends on the person's immune response; poor cell-mediated immunity leads to the lepromatous form with many bacteria
    \item \textbf{Nerve damage:} the bacterium prefers cooler areas of the body and invades peripheral nerves, causing loss of sensation
    \item \textbf{Skin involvement:} skin lesions vary from a few pale patches (tuberculoid form) to widespread nodules (lepromatous form)
    \item \textbf{Disability:} nerve damage causes loss of feeling, weakness, and, without care, injury, ulcers, and deformity
    \item \textbf{Not highly contagious:} most people are naturally immune
\end{itemize}

\section*{Clinical Features}
Symptoms develop slowly and vary with the form of the disease.

\begin{itemize}
    \item Pale or reddish skin patches with reduced sensation
    \item Numbness, tingling, or loss of feeling in the hands, feet, or face
    \item Weakness of the hands or feet
    \item Thickened or painful nerves
    \item Skin nodules or thickening of the skin in the lepromatous form
    \item Loss of eyebrows and eyelashes in advanced disease
    \item Nasal congestion and nosebleeds
    \item Painless injuries, burns, or ulcers because of lost sensation
    \item Eye problems, including dryness and loss of blinking, from facial nerve damage
\end{itemize}

\section*{Differential Diagnosis}
Other skin and nerve conditions can resemble Hansen's disease.

\begin{itemize}
    \item Other causes of skin patches, including tinea, eczema, and vitiligo
    \item Other causes of peripheral neuropathy, including diabetes and vitamin deficiency
    \item Sarcoidosis and other granulomatous skin diseases
    \item Fungal infections of the skin
    \item Cutaneous leishmaniasis and other tropical infections
    \item Lupus and other autoimmune skin conditions
    \item Peripheral nerve entrapment
\end{itemize}

\section*{When to Seek Medical Care}
See a doctor for skin patches with reduced sensation, persistent numbness or weakness, or skin lesions that are unusual and not responding to treatment, especially if you have lived in or travelled to a country where Hansen's disease occurs.

\textbf{Call 000 (or local emergency services) for:}
\begin{itemize}
    \item This condition is not an emergency; however, seek urgent care for new weakness or loss of sensation that is sudden or progressing rapidly
    \item Signs of infection in areas with lost sensation, including spreading redness, swelling, and fever
    \item Severe pain or injury in a numb limb
\end{itemize}

\section*{Diagnosis and Investigations}
Hansen's disease is diagnosed by clinical examination and laboratory tests.

\begin{itemize}
    \item \textbf{Clinical examination:} assessment of skin lesions and sensation, and examination of peripheral nerves
    \item \textbf{Skin smear:} microscopic examination of skin slit smears for the bacterium
    \item \textbf{Skin biopsy:} examined under a microscope for characteristic features
    \item \textbf{Nerve examination:} testing sensation, strength, and reflexes
    \item \textbf{PCR testing:} molecular confirmation of the organism
    \item \textbf{Classification:} cases are classified as paucibacillary or multibacillary to guide treatment
    \item \textbf{Notification:} Hansen's disease is notifiable in Australia
\end{itemize}

\section*{Management}
Multidrug therapy cures Hansen's disease and prevents disability.

\begin{itemize}
    \item \textbf{Multidrug therapy:} a combination of antibiotics, including dapsone and rifampicin, taken for six to twelve months or longer depending on the form
    \item \textbf{Free treatment:} treatment is provided free through the World Health Organization
    \item \textbf{Prevention of disability:} protecting numb areas from injury, footwear care, and eye care
    \item \textbf{Physiotherapy:} exercises to maintain strength and prevent contractures
    \item \textbf{Management of reactions:} immune reactions to the bacteria can cause pain and swelling and are treated with steroids
    \item \textbf{Specialist care:} management in a specialist infectious diseases or leprosy service
    \item \textbf{Contact tracing:} close contacts are examined and may be offered preventive treatment
    \item \textbf{Rehabilitation:} support for people with established disability
\end{itemize}

\section*{Complications}
\begin{itemize}
    \item Permanent nerve damage causing loss of sensation and weakness
    \item Painless injuries, burns, and ulcers in numb areas
    \item Deformity of the hands, feet, and face if untreated
    \item Eye damage and blindness from facial nerve involvement
    \item Immune reactions during or after treatment
    \item Stigma and psychological distress
    \item Social and economic effects of disability
\end{itemize}

\section*{Prognosis and Outcomes}
With early diagnosis and treatment, Hansen's disease is curable, and disability is largely prevented. Multidrug therapy kills the bacteria, stops transmission, and cures the infection within months to a year or two. Nerve damage that has already occurred cannot be reversed, which is why early diagnosis is so important. With treatment and rehabilitation, most people with Hansen's disease live full and productive lives, and the risk of transmission to others is eliminated once treatment begins.

\section*{Prevention and Self-Care}
\begin{itemize}
    \item Seek early assessment of skin patches with reduced sensation or unusual skin lesions
    \item If diagnosed, complete the full course of treatment
    \item Protect numb areas from injury: check feet daily, wear protective footwear, and avoid burns
    \item Attend physiotherapy and follow exercise programs
    \item Protect the eyes if blinking is reduced, with artificial tears and care
    \item Practise good skin and wound care
    \item Encourage close contacts to be examined
    \item Address stigma by understanding that the disease is curable and not highly contagious
    \item Seek support for the psychological effects of the condition
\end{itemize}

\section*{Sources and Further Reading}
The following reputable sources provide further information for healthcare students and patients seeking reliable, current advice about Hansen's disease.

\begin{itemize}
    \item Australian Government Department of Health and Aged Care. \textit{Leprosy (Hansen's disease) fact sheet.} https://www.health.gov.au/
    \item Healthdirect Australia. \textit{Hansen's disease (leprosy).} https://www.healthdirect.gov.au/
    \item NSW Health. \textit{Leprosy fact sheet.} https://www.health.nsw.gov.au/
    \item World Health Organization. \textit{Leprosy fact sheet.} https://www.who.int/
    \item Centers for Disease Control and Prevention. \textit{Hansen's disease (leprosy).} https://www.cdc.gov/
    \item International Federation of Anti-Leprosy Associations. \textit{Information about leprosy.} https://ilepfederation.org/
\end{itemize}
